\section{Conclusie}
\label{sec:conclusie}

In dit project is met behulp van sensorfusie onderzocht of de aanwezige IMU de 
positiebepaling van de ProRTK tracker robuuster en betrouwbaarder kan maken. Hiervoor is een Extended Kalman Filter ontworpen en gevalideerd, getoetst aan de criteria uit hoofdstuk~\ref{sec:specificaties}. Onder normale omstandigheden volgt het filter de RTK-positie binnen 1\,cm en blijft daarmee binnen het criterium van 2\,cm, zonder de bestaande nauwkeurigheid aan te tasten. \\

Een grote winst van het systeem blijkt tijdens verstoringen. Foutieve sprongen door multipath worden herkend en verworpen, en bij uitval van twee seconden blijft de afwijking bij een stabiele snelheid binnen de gestelde 1\,m. Bij sterke versnelling of vertraging wordt deze grens echter overschreden, met afwijkingen tot 3,9\,m. Hiermee kan de hoofdvraag worden beantwoord: sensorfusie met de aanwezige IMU verbetert de robuustheid en betrouwbaarheid merkbaar en vangt de belangrijkste tekortkomingen van een losstaand RTK systeem grotendeels op. \\

Van de onderzochte aanvullende sensoren is de aanwezige IMU binnen de hardware randvoorwaarde het meest geschikt. De magnetometer is de meest kansrijke toevoeging voor de toekomst, maar vereist een 9-assige IMU. Voor toepassing in paardenraces is aanvullende validatie nodig, omdat de huidige tests de schokken en dynamiek van een galopperend paard niet volledig dekken. Hierom wordt aanbevolen het systeem te valideren met een rijdend paard vergelijkbaar met de validatietest tijdens dit project.

\newpage

\section{Discussie}
\label{sec:discussie}

Kijkend naar het gehele systeem presteert deze goed onder de omschreven omstandigheden. Echter bevat het een aantal beperkingen. \\

De prestaties van het systeem tijdens een signaaluitval zijn niet altijd consistent. Bij een stabiele snelheid en een korte uitval overbrugt het systeem het gat nauwkeurig, met afwijkingen onder de meter. Op momenten van sterke versnelling of vertraging loopt de afwijking echter op tot enkele meters, met een maximum van ongeveer 3,9\,m. Het systeem werkt dus in de meeste gevallen goed, maar bij abrupte veranderingen in snelheid kan het de positie flink over- of onderschatten. \\

Een tweede beperking is het gedrag bij achteruitrijden. De u-blox levert de snelheid altijd als een positieve waarde, terwijl het model uitgaat van een voorwaartse beweging in de kijkrichting. Bij achteruitrijden komt de gemeten beweging daardoor niet overeen met de werkelijkheid en wordt de positie verkeerd ingeschat. Voor de paardensport, waarin uitsluitend voorwaarts wordt bewogen, vormt dit geen probleem, maar voor toepassing bij autosport is het een aandachtspunt. \\

Tot slot is de validatie uitgevoerd met loop- en autotests onder gecontroleerde 
omstandigheden. Deze tonen aan dat het filter correct functioneert, maar dekken niet de dynamiek van een galopperend paard. De tracker kan volgens de specificaties schokken tot ongeveer 2,5\,$g$ ervaren, en de invloed daarvan op de metingen is in de huidige tests niet onderzocht. De resultaten gelden daarom als bewijs dat het systeem onder de geteste omstandigheden werkt, maar nog niet als bewijs voor de uiteindelijke toepassing in de paardensport.


\newpage